% generated from Paper 7/anccft7.tex
\renewcommand{\ind}{\operatorname{ind}_{\mathbf1}}
\chapter{Algorithmic non-commutative class field theory, VII: the snake factorization of the transfer and the evaluated characteristic form}\label{chap:VII}
\chaptermark{VII}
\begin{chapterabstract}
Problem 5.1 of [Ch.~\ref{chap:V}] asked for a characteristic formalism on the degeneracy tower
satisfying three axioms: an arc-density term recovering $\mu=n(q+1)/q$ (A1), an
additive Euler term $-\chi k$ with $\chi$ the rank of the degeneracy $H^0$ (A2), and
reduction to the commutative characteristic ideal of [Ch.~\ref{chap:IV}] on abelian sub-towers
(A3). This paper solves the problem in evaluated form, by a single application of
the snake lemma, and the mechanism turns out to be simpler than anything the series
anticipated: because the level-$(k{+}1)$ Hecke operator factors through the tail
norm, $A_{k+1}=\etp\ta$ {\rm[Ch.~\ref{chap:V}, Prop.~1.2]}, it \emph{vanishes identically} on
$\ker\ta$, so the Laplacian restricted to the norm kernel is exactly $qI$ and the
``arc determinant'' is exactly $q^{\,n_k(q-1)}$ --- no nilpotency estimate, no unit
ambiguity. The snake lemma applied to
$0\to\ker\nu\to\Zp^{N_{k+1}}\to\Zp^{N_k}\to0$ with Laplacian verticals then yields,
for every exactly intertwining, surjective, degree-preserving norm $\nu$,
\[
\ord_p\Tor_{k+1}-\ord_p\Tor_k
=\ord_p\bigl|\coker(\Delta_{k+1}|_{\ker\nu})\bigr|-\ord_p\ind(\nu),
\]
where $\ind(\nu)$ is the factor by which $\nu$ scales the Eisenstein line. For
$\nu=\ta$ this gives $n_k(q-1)-1$: axiom (A1) with $\mathrm{Det_{arc}}
=q^{\,n_k(q-1)}$ and axiom (A2) with $\mathrm{Det_{Euler}}=\ind(\ta)=q$, one factor
per level, $\chi=1$. For $\nu=\pi_*$ on an abelian voltage tower the same snake
gives $\ord_\ell\prod_{\text{new }\zeta}\det D(\zeta)-1$: axiom (A3), the character
mechanism of [Ch.~\ref{chap:II}, \S5] and [Ch.~\ref{chap:IV}], verified to the integer
($\det(\Delta|_{\ker\pi})=48$ and $784$ against the character products). One
formula, two specializations, all three axioms now theorems. What is not delivered
--- a lift of this evaluated form to a single class in $K_1$ of a localized
completed degeneracy algebra --- is stated as the precise residue of Problem 5.1,
and, consistently with [Ch.~\ref{chap:III}] and [Ch.~\ref{chap:V}], we decline to axiomatize a localization theory
that does not exist.
\end{chapterabstract}

\section*{Status of the results}

All numbered statements are proved; the proofs use only results of Papers I--VI and
the snake lemma. Every identity was additionally machine-verified exactly
(Table~\\ref{VII:tab:battery}): Iwahori transitions $1\to2$, $2\to3$ over $K_4$, and
abelian transitions $0\to1$, $1\to2$ of the voltage tower $(1,1,1)$, $\ell=2$.
Remark~\\ref{VII:rem:residue} states what remains open of [Ch.~\ref{chap:V}, Prob.~5.1];
Remark~\\ref{VII:rem:corrections} records two corrections to the task specification that
prompted this paper. There are no external citations.

\section{The kernel of the norm, and the vanishing lemma}\label{VII:sec:kernel}

Setting as in [Ch.~\ref{chap:V}]: $Y_k$ the Iwahori path tower over a $(q+1)$-regular base
($n_k=n(q+1)q^{k-1}$ for $k\ge1$), $\ta$ the tail norm, $A_k$ the level Hecke
(adjacency) operator, $\Delta_k=qI-A_k$ for $k\ge1$, sandpile torsion
$\Tor_k$.

\begin{lemma}[Vanishing on the norm kernel]\label{VII:lem:vanish}
For $k\ge1$ let $K_k:=\ker_{\Zp}(\ta)\subseteq\Zp^{\,n_{k+1}}$. Then:
\begin{enumerate}
\item[(i)] $K_k$ is a saturated direct summand of rank $n_{k+1}-n_k=n_k(q-1)$
(kernels of integer matrices are saturated);
\item[(ii)] $A_{k+1}$ vanishes on $K_k$: since $A_{k+1}=\etp\,\ta$
{\rm[Ch.~\ref{chap:V}, Prop.~1.2]}, $A_{k+1}x=\etp(\ta x)=0$ for $x\in K_k$; in particular $K_k$
is invariant and
\[
\Delta_{k+1}\big|_{K_k}=q\,\mathrm{id}_{K_k},
\qquad
\bigl|\coker(\Delta_{k+1}|_{K_k})\bigr|=q^{\,n_k(q-1)}\ \text{exactly.}
\]
\end{enumerate}
\end{lemma}

\begin{proof}
(i) is the saturation principle of [Ch.~\ref{chap:II}] (if $ka\in\ker$, $k\neq0$, then
$a\in\ker$); the rank is surjectivity of $\ta$. (ii) is the displayed one-line
computation; then $\Delta_{k+1}=qI-A_{k+1}$ restricts to $q\,\mathrm{id}$, whose
cokernel on a rank-$n_k(q-1)$ lattice has order exactly $q^{\,n_k(q-1)}$. Verified:
the machine search for the nilpotency index of $A_{k+1}|_{K_k}$ returned $s=1$ at
both computed transitions, and $\det(\Delta_{k+1}|_{K_k})=2^{12},2^{24}$ on the
nose.
\end{proof}

Lemma~\\ref{VII:lem:vanish}(ii) is the structural surprise of this paper: the series
expected, at best, a nilpotency estimate through the characteristic-polynomial ratio
of [Ch.~\ref{chap:III}, eq.~(2.1)]; instead the incidence calculus of [Ch.~\ref{chap:V}] kills the operator
outright, because the Hecke operator of the upper level literally factors through
the norm to the lower one.

\section{The snake factorization}\label{VII:sec:snake}

\begin{theorem}[Snake factorization of the transfer]\label{VII:thm:snake}
Let $\nu\colon\Zp^{N_{k+1}}\to\Zp^{N_k}$ be surjective and degree-preserving, with
$\nu\Delta_{k+1}=\Delta_k\nu$, and suppose the Eisenstein lines satisfy
$\nu(\mathbf1)=c\,\mathbf1$ for an integer $c=:\ind(\nu)\neq0$. Then $K:=\ker\nu$ is
a saturated invariant summand, the snake lemma applied to
\[
0\longrightarrow K\longrightarrow \Zp^{N_{k+1}}
\xrightarrow{\ \nu\ }\Zp^{N_k}\longrightarrow0
\]
with verticals $\Delta_{k+1}|_K$, $\Delta_{k+1}$, $\Delta_k$ yields the exact
sequence
\[
0\to\Zp\xrightarrow{\ \cdot c\ }\Zp\xrightarrow{\ \delta\ }
\coker\bigl(\Delta_{k+1}|_K\bigr)\to\coker\Delta_{k+1}\to\coker\Delta_k\to0 ,
\]
and hence, whenever $\coker(\Delta_{k+1}|_K)$ is finite,
\begin{equation}\label{VII:eq:master}
\ord_p\Tor_{k+1}-\ord_p\Tor_k
\;=\;\ord_p\bigl|\coker(\Delta_{k+1}|_{K})\bigr|\;-\;\ord_p\,c .
\end{equation}
\end{theorem}

\begin{proof}
Invariance of $K$ is the intertwining; saturation as in
Lemma~\\ref{VII:lem:vanish}(i). The kernels of the verticals are:
$\ker(\Delta_{k+1}|_K)=K\cap\Zp\mathbf1=0$ (the constants are not in $\ker\nu$
since $\nu\mathbf1=c\mathbf1\neq0$), $\ker\Delta_{k+1}=\Zp\mathbf1$,
$\ker\Delta_k=\Zp\mathbf1$, and the induced map between them is
$\mathbf1\mapsto c\mathbf1$. The snake lemma gives the displayed sequence; the free
ranks of the two cokernels match through $\nu$ (degree preservation, as in
[Ch.~\ref{chap:IV}, Lem.~1.3]), so taking orders of the finite parts yields \\eqref{VII:eq:master}.
\end{proof}

\begin{corollary}[Iwahori specialization: axioms (A1) and (A2)]\label{VII:cor:iwahori}
For $\nu=\ta$, $k\ge1$: $c=q$ $($[V], $\ta\mathbf1=q\mathbf1)$ and
Lemma~\\ref{VII:lem:vanish} give
\[
\ord_p\Tor_{k+1}-\ord_p\Tor_k=n_k(q-1)-1 ,
\]
recovering {\rm[Ch.~\ref{chap:III}, Thm.~2.1]} homologically. Summation yields the pseudo-Iwasawa
law with $\mu=n(q+1)/q$ carried entirely by
$\mathrm{Det_{arc}}:=\det(\Delta_{k+1}|_{\ker\ta})=q^{\,n_k(q-1)}$ (axiom A1), and
the Euler term $-\chi k$ carried entirely by
$\mathrm{Det_{Euler}}:=\ind(\ta)=q$, one factor per level, with
$\chi=1=\rk\ker(\tau^*-\eta^*)$ by {\rm[Ch.~\ref{chap:V}, Thm.~2.1]} (axiom A2): the $-1$ is,
finally and precisely, the Eisenstein line passing through the connecting map of the
snake.
\end{corollary}

\begin{corollary}[Abelian specialization: axiom (A3)]\label{VII:cor:abelian}
For an abelian $\Z/\ell^{k}$ voltage tower with covering pushforward $\nu=\pi_*$
$($exactly intertwining by {\rm[Ch.~\ref{chap:IV}, Lem.~1.2]}$)$: $c=\ell$, and under the
block-character decomposition of {\rm[Ch.~\ref{chap:IV}, Lem.~1.1]} the kernel $\ker\pi_*$ is the
sum of the \emph{new} character blocks, on which $\Delta$ acts with determinant
$\prod_{\text{new }\zeta}\det D(\zeta)$; hence \\eqref{VII:eq:master} reads
\[
\ord_\ell\Tor_{k+1}-\ord_\ell\Tor_k
=\ord_\ell\!\!\prod_{\text{new }\zeta}\!\det D(\zeta)\;-\;1 ,
\]
which is precisely the character/Newton mechanism of {\rm[Ch.~\ref{chap:II}, \S5]} and
{\rm[Ch.~\ref{chap:IV}, Cor.~2.3]}. Verified exactly: $\det(\Delta|_{\ker\pi})=48=\det D(-1)$ at
$0\to1$ and $784=|\det D(i)\det D(-i)|$ at $1\to2$, increments $4-1=3$ in both,
matching {\rm[Ch.~\ref{chap:II}, Table~2]}.
\end{corollary}

One snake, two kernels: on the Iwahori side the kernel is killed by the Hecke
operator and contributes the pure arc power $q^{\,n_k(q-1)}$; on the abelian side
the kernel is the new-character space and contributes the $L$-value block
$\prod_{\text{new}}\det D(\zeta)$. The Eisenstein index $c$ is $q$ or $\ell$
respectively --- one factor per level in either world, which is the whole content of
$\lambda=-\chi$. The three axioms of [Ch.~\ref{chap:V}, \S5] are now theorems with these referents.

\section{Scope, residue, corrections}\label{VII:sec:scope}

\begin{remark}[The residue of Problem 5.1]\label{VII:rem:residue}
Theorem~\\ref{VII:thm:snake} solves [Ch.~\ref{chap:V}, Prob.~5.1] \emph{in evaluated form}: the
characteristic form is the level-wise pair
$\mathrm{Char}_k:=\bigl(\det(\Delta_{k+1}|_{\ker\nu}),\ \ind(\nu)\bigr)$, and
(A1)--(A3) hold for it as theorems. What is not constructed is a single element of
$K_1$ of a localized completion of the degeneracy algebra --- a Dieudonn\'e-type
determinant --- whose evaluations reproduce $\mathrm{Char}_k$. Defining such an
object would require an Ore localization theory for the non-normal pro-ring of
[Ch.~\ref{chap:V}, Prop.~4.1], which does not exist; consistently with [Ch.~\ref{chap:III}, Rem.~2.6] and
[Ch.~\ref{chap:V}, \S5], we decline to axiomatize it, and we record the lift as the precise open
residue. Everything quantitative that a lift could evaluate to is already proved
here.
\end{remark}

\begin{remark}[Two corrections to the task specification]\label{VII:rem:corrections}
The specification prompting this paper asserted an exact sequence
$0\to X_{k+1}\xrightarrow{\partial}X_k\to0$ with $\partial=\ta-\et$; that map is far
from injective ($\rk\ker\partial=n_{k+1}-n_k+1$ by {\rm[Ch.~\ref{chap:V}, Thm.~2.1(ii)]}), and the
exact sequence that carries the theory is the norm sequence of
Theorem~\\ref{VII:thm:snake}. It also placed the Eisenstein class in
$K_0(\mathcal H_k)$; the Eisenstein \emph{projector} $\mathbf1\mathbf1^{t}/N$ is not
integral and does not lie in the algebra, and the class that exists is the
module-level line $\Zp\mathbf1=\ker\Delta$, entering the theory through the
connecting map --- which is where Corollary~\\ref{VII:cor:iwahori} puts it.
\end{remark}

\begin{remark}[Boundary level]\label{VII:rem:boundary}
At $k=0$ the tail map fails to intertwine by exactly the degeneracy boundary
$($[Ch.~\ref{chap:V}, Rem.~3.3]$)$, $\ker S$ is not invariant, and Theorem~\\ref{VII:thm:snake} does not
apply; the Eisenstein index there is $q+1$, not $q$
$(S\mathbf1=(q+1)\mathbf1$, verified$)$. The transfer $0\to1$ enters the law as the
initial value, as everywhere in the series.
\end{remark}

\section{The verification battery}\label{VII:sec:battery}

\begin{table}[ht]
\centering\small
\renewcommand{\arraystretch}{1.2}
\begin{tabular}{lccccc}
\toprule
transition & $\rk\ker\nu$ & invariant & $A_{k+1}|_{\ker}$ &
$\det(\Delta|_{\ker\nu})$ & $\ind(\nu)$\\
\midrule
Iwahori $1\to2$ & $12$ & yes & $=0$ ($s{=}1$) & $2^{12}$ exactly & $2$\\
Iwahori $2\to3$ & $24$ & yes & $=0$ ($s{=}1$) & $2^{24}$ exactly & $2$\\
abelian $0\to1$ & $4$ & yes & --- & $48=\det D(-1)$ & $2$\\
abelian $1\to2$ & $8$ & yes & --- & $784=|\det D(i)\det D(-i)|$ & $2$\\
boundary $0\to1$ & --- & no & --- & --- & $3=q+1$\\
\bottomrule
\end{tabular}
\medskip
\caption{Exact verification (\texttt{iwahori\_char.py}). Increments
\\eqref{VII:eq:master}: $12-1=11$, $24-1=23$ (Iwahori, matching [Ch.~\ref{chap:II}, Rem.~5.2]);
$4-1=3$, $4-1=3$ (abelian $2$-adic orders, matching [Ch.~\ref{chap:II}, Table~2]). The vanishing
$A_{k+1}|_{\ker\ta}=0$ was found by the machine ($s=1$) before its one-line proof
was noticed.}
\label{VII:tab:battery}
\end{table}

\section*{Acknowledgement of provenance and caveats}

Unrefereed preprint. Every numbered statement is proved from Papers I--VI and the
snake lemma; the battery of Table~\\ref{VII:tab:battery} confirms each ingredient
exactly. No external sources are cited and none are load-bearing.
Remark~\\ref{VII:rem:residue} states the open residue of [Ch.~\ref{chap:V}, Prob.~5.1];
Remark~\\ref{VII:rem:corrections} corrects the specification that prompted the paper.
